% =====================================================================
% DRAFT v0 (01.11.2026) -- k-bonacci neighborhoods for the PFSP
% Target: ICAPS 2027 (fallback IJCAI 2027). Both use double-blind review.
% Layout: plain article class for drafting; switch to the official
% AAAI/ICAPS (or IJCAI) template once the CFP is out.
% Self-citations are written in the third person (double-blind).
% Numbers in Table 2 come from the PILOT (scripts/run_kbonacci_pilot.py);
% they are placeholders for the full Taillard campaign.
% =====================================================================
\documentclass[10pt,twocolumn]{article}
\usepackage[margin=0.75in]{geometry}
\usepackage[T1]{fontenc}
\usepackage{amsmath,amssymb,amsthm}
\usepackage{booktabs}
\usepackage{algorithm}
\usepackage{algpseudocode}
\usepackage{xcolor}
\usepackage{hyperref}

\newtheorem{proposition}{Proposition}
\newtheorem{definition}{Definition}
\newcommand{\todo}[1]{\textcolor{red}{[TODO: #1]}}
\newcommand{\N}{\mathcal{N}}

\title{From Fibonacci to Dynasearch: the $k$-bonacci Family of\\
Compound-Swap Neighborhoods for the Permutation Flow Shop}
\author{Anonymous submission}
\date{}

\begin{document}
\maketitle

\begin{abstract}
Compound-move neighborhoods search an exponential set of schedules in
polynomial time by dynamic programming, but the two best-known members
of the family for the permutation flow shop sit at opposite extremes:
the Fibonacci neighborhood combines only adjacent swaps and is cheap,
while Dynasearch combines swaps of any distance and costs $O(mn^3)$ per
evaluation. We introduce the \emph{$k$-bonacci} family that interpolates
between them: a move is a set of endpoint swaps on disjoint intervals
whose swap distance is at most $k-1$. The family is counted exactly by
the $k$-bonacci numbers, so its size grows as $\lambda_k^n$ with
$\lambda_2=1.618$ (Fibonacci), $\lambda_3=1.839$ (tribonacci) and
$\lambda_k\to 2$ (Dynasearch). We give a positional dynamic program
that finds the best move in $O(nk)$ time after $O(mnk^2)$ delta
evaluation, combine it with the block-property accelerator, and study
the parameter $k$ empirically: solution quality, cost per iteration,
structure of the selected moves, and the error of the additivity
assumption that every compound-move DP makes in the flow shop.
\todo{one sentence with the headline result of the full campaign}
\end{abstract}

% ---------------------------------------------------------------------
\section{Introduction}

The permutation flow shop scheduling problem (PFSP) with makespan
criterion, $F|prmu|C_{\max}$, is NP-hard for $m\geq 3$
machines~\cite{garey1976} and remains a standard testbed for
local search~\cite{taillard1993,ruiz2007}. Two design decisions dominate
the performance of a local search for the PFSP: the move structure, and
the critical-path (block) properties that prune it~\cite{nowicki1996,grabowski2004}.

Dynasearch~\cite{congram2002} showed that an exponentially large
neighborhood of \emph{independent} pairwise interchanges can be searched
exactly by dynamic programming, turning a compound move into the unit
of search; the idea belongs to the broader class of very large-scale
neighborhood search (VLSN)~\cite{ahuja2002}. Recent work introduced two
further compound neighborhoods for the PFSP whose sizes are counted by
classical integer sequences: the Fibonacci neighborhood of disjoint
adjacent swaps~\cite{fibonacci_own,bozejko2026sr} and the Motzkin
neighborhood of non-crossing, possibly nested swaps~\cite{motzkin_own}.

These neighborhoods differ by orders of magnitude in size and cost, yet
they are always compared as isolated points. This paper fills the space
between Fibonacci and Dynasearch with a one-parameter family. Our
contributions are:
\begin{enumerate}
  \item the \textbf{$k$-bonacci neighborhood} $\N_k$, with the
        chain $\N_{\mathrm{adj}}\subset\N_2=\N_{\mathrm{fib}}\subset\N_3
        \subset\dots\subset\N_n=\N_{\mathrm{dyn}}$, and an exact count
        of $|\N_k|$ by the $k$-bonacci recurrence (Prop.~\ref{prop:count});
  \item a \textbf{positional DP} that selects the best compound move in
        $O(nk)$ time, compatible with the block-property accelerator;
        for $k=n$ it is also a faster, non-recursive Dynasearch;
  \item an \textbf{empirical study of $k$} on the Taillard
        benchmark~\cite{taillard1993}: quality versus $k$, cost per
        iteration versus $k$, move anatomy, and the
        \textbf{additivity error} of compound-move DP in the PFSP;
  \item \todo{RQ5 adaptive $k$ (IJCAI variant) or comparison with
        Iterated Greedy~\cite{ruiz2007}}.
\end{enumerate}

% ---------------------------------------------------------------------
\section{Related Work}

\paragraph{Compound moves.} Dynasearch was introduced for single-machine
total weighted tardiness~\cite{congram2002}; Ergun and
Orlin~\cite{ergun2006} reduced its swap neighborhood search to
$O(n^2)$, and Grosso et al.~\cite{grosso2004} added insertions.
Bo\.zejko and Wodecki analyzed swap multimoves for the
PFSP~\cite{bozejko2007}. In single-machine total weighted tardiness the
cost change of independent swaps is additive, which makes the DP exact;
in the PFSP it is not (Section~\ref{sec:additivity}).
\paragraph{Sequence-counted neighborhoods.} \todo{Fibonacci (own work,
third person), Sci.\ Rep.\ 2026~\cite{bozejko2026sr}, IFAC 2025 on
non-crossing intervals -- verify that neither contains the $k$-bonacci
hierarchy}. Motzkin~\cite{motzkin_own}.
\paragraph{Parameterized local search.} Balzereit et
al.~\cite{balzereit2024} study the radius--runtime trade-off for
distance-bounded local search in single-machine scheduling from a
parameterized-complexity viewpoint; we study the same kind of trade-off
empirically for a compound neighborhood.
\paragraph{PFSP state of the art.} NEH~\cite{nawaz1983}, tabu
search~\cite{nowicki1996,grabowski2004}, Iterated
Greedy~\cite{ruiz2007}. \todo{2--3 recent papers}

% ---------------------------------------------------------------------
\section{Problem and Notation}

$n$ jobs are processed on machines $1,\dots,m$ in the same order
$\pi=(\pi(1),\dots,\pi(n))$; $p_{r,j}$ is the processing time of job $j$
on machine $r$. Completion times satisfy
$C_{r,t}=\max(C_{r-1,t},C_{r,t-1})+p_{r,\pi(t)}$ and
$C_{\max}(\pi)=C_{m,n}$. Head and tail matrices $H$ and $T$ (forward and
backward longest paths) give the makespan of any modification of a
segment $[i,j]$ by recomputing only that segment:
$C_{\max}=\max_r\bigl(H'_{r,j}+T_{r,j+1}\bigr)$, at cost $O(m(j-i+1))$.

% ---------------------------------------------------------------------
\section{The $k$-bonacci Neighborhood}

\begin{definition}
Let $k\geq 2$. An \emph{elementary move} $s(i,j)$, $1\leq j-i\leq k-1$,
exchanges $\pi(i)$ and $\pi(j)$ and leaves $\pi(i+1),\dots,\pi(j-1)$ in
place. A \emph{compound move} is a set of elementary moves whose
intervals $[i,j]$ are pairwise disjoint. $\N_k(\pi)$ is the set of
permutations reachable by one compound move.
\end{definition}

Special cases: $k=2$ allows swap distance $1$ only (Fibonacci), $k=3$
distances $1,2$ (tribonacci), $k=n$ any distance (Dynasearch).

\begin{proposition}\label{prop:count}
$|\N_k(\pi)|=a_n$, where $a_0=1$, $a_p=\sum_{L=1}^{\min(k,p)}a_{p-L}$.
Hence $|\N_2|=F_{n+1}$, $|\N_n|=2^{n-1}$, and
$|\N_k|=\Theta(\lambda_k^n)$, where $\lambda_k$ is the largest real
root of $x^k=x^{k-1}+\dots+x+1$.
\end{proposition}
\begin{proof}[Sketch]
A compound move is a tiling of positions $1..n$ by tiles of length $1$
(untouched position) and $L\in\{2,\dots,k\}$ (swap the tile ends);
conditioning on the first tile gives the recurrence. Distinct tilings
give distinct permutations: the tile $[i,j]$ is the smallest interval
starting at $i$ mapped onto itself, so the tiling can be read back from
the permutation. \todo{full proof}
\end{proof}

\begin{table}[t]
\centering
\caption{Size of $\N_k$ for $n=10$ and growth constant $\lambda_k$.}
\label{tab:count}
\begin{tabular}{@{}lrrrrrr@{}}
\toprule
$k$ & 2 & 3 & 4 & 5 & 6 & $n$ \\
\midrule
$|\N_k|$, $n=10$ & 89 & 274 & 401 & 464 & 492 & 512 \\
$\lambda_k$ & 1.618 & 1.839 & 1.928 & 1.966 & 1.984 & 2 \\
\bottomrule
\end{tabular}
\end{table}

\paragraph{Evaluation.} For every candidate $(i,j)$ we compute
$\delta_{ij}=C_{\max}(\pi\circ s(i,j))-C_{\max}(\pi)$ with heads and
tails; there are fewer than $n(k-1)$ candidates, each costing
$O(m(j-i))$, so all deltas cost $O(mnk^2)$, against $O(mn^3)$ for
Dynasearch.
\paragraph{Block-property accelerator.} A swap $(i,j)$ can decrease
$C_{\max}$ only if a block boundary of the critical path lies in
$[i,j]$~\cite{smutnicki1998,nowicki1996}. Boundaries are found in
$O(mn)$ from $H$, and a running flag keeps the filter at $O(1)$ per
candidate.
\paragraph{Positional DP.} Let $B[p]$ be the minimum total delta of a
compound move restricted to positions $1..p$. Then
\begin{equation}
B[p]=\min\Bigl(B[p-1],\ \min_{1\le d\le k-1} B[p-d-1]+\delta_{p-d,p}\Bigr),
\end{equation}
with $B[0]=0$, which takes $O(nk)$ time and memory. When no move improves
the predicted makespan, the best single swap is applied, as in
Dynasearch.

% ---------------------------------------------------------------------
\section{Additivity}\label{sec:additivity}
The DP optimizes $\sum\delta_{ij}$, but in the PFSP the makespan change
of a compound move is not the sum of the individual changes: two swaps
on disjoint intervals interact through the critical path. Any DP-based
compound neighborhood (Fibonacci, Dynasearch, Motzkin) therefore
optimizes a surrogate. Larger $k$ admits more and longer swaps per move,
so the surrogate error should grow with $k$; Section~\ref{sec:exp}
measures it.

% ---------------------------------------------------------------------
\section{Experiments}\label{sec:exp}

\paragraph{Setup (pilot).} Taillard instances $20\times5$, $20\times10$
and $50\times10$ (instances 1--4 of each), 3 seeds, ILS with tabu list
and elite-pool perturbation~\cite{windowed_own}, time limit
$1000$\,ms, Python~3.13, 4 runs in parallel. RPD is measured with respect
to the upper bounds in the Taillard files. \todo{full campaign:
12 classes $\times$ 10 instances, 5 seeds, budgets 100--10\,000 ms,
best-known upper bounds, Friedman + Nemenyi, IG baseline}

\begin{table}[t]
\centering
\caption{PILOT: mean RPD [\%] to the upper bound / mean ILS iterations
in 1000\,ms (4 instances $\times$ 3 seeds per class).}
\label{tab:pilot}
\small
\begin{tabular}{@{}lccc@{}}
\toprule
Neighborhood & $20\times5$ & $20\times10$ & $50\times10$ \\
\midrule
Adjacent         & 1.59 / 4578 & 3.26 / 2246 & 11.73 / 403 \\
Fibonacci ($k{=}2$) & 1.31 / 12400 & 2.46 / 7129 & 8.37 / 3329 \\
tribonacci ($k{=}3$) & 0.92 / 10177 & 1.84 / 5007 & 5.94 / 2481 \\
$k=4$            & 0.89 / 7897 & 1.43 / 3618 & 5.51 / 1966 \\
$k=6$            & 0.74 / 5397 & \textbf{1.09} / 2375 & 4.85 / 1123 \\
$k=8$            & 0.72 / 3893 & 1.35 / 1580 & 4.66 / 809 \\
$k=n$            & \textbf{0.55} / 1539 & \textbf{0.99} / 610 & \textbf{3.62} / 63 \\
Dynasearch (recursive) & 0.94 / 655 & 1.49 / 384 & 5.30 / 23 \\
\bottomrule
\end{tabular}
\end{table}

\paragraph{Quality versus $k$ (pilot).} Table~\ref{tab:pilot} shows a
monotone improvement from $k=2$ to $k=6$ on all three classes, and at
$1000$\,ms the full $k=n$ is still best, although its iteration count
collapses from $3329$ to $63$ at $n=50$. The positional DP makes $k=n$
$2$--$3\times$ faster than the recursive Dynasearch implementation with
the same move set, which alone improves RPD from $5.30\%$ to $3.62\%$ at
$50\times10$. Whether an interior optimum $k^*$ appears for $n\geq 100$
or for shorter budgets is the main question for the full campaign.
\todo{Fig.: RPD vs $k$ per $n$ and budget; iterations vs $k$}

\paragraph{Additivity error (pilot).} On random permutations the predicted
delta is exact in 42\% of evaluations for $k=2$, 16--20\% for $k=3,4$
and only 10--12\% for $k\geq 6$. The predicted improvement overestimates
the real one by 10\% at $k=2$, 26\% at $k=3,4$, 41\% at $k=6$ and about
50\% at $k\geq 8$.
The real improvement nevertheless keeps growing with $k$.
\todo{same measurement along ILS trajectories (near local optima)}

\paragraph{Move anatomy.} A selected move contains 2--5 swaps on average,
almost independently of $k$, while the mean swap distance grows from
$1$ ($k=2$) to $4.4$--$4.6$ ($k=8$) and $5$--$9$ ($k=n$): larger $k$ is
used to make longer swaps, not more of them. \todo{histograms}

% ---------------------------------------------------------------------
\section{Conclusions}
\todo{after the full campaign: recommended $k$, cost model,
additivity-aware variant (re-evaluation of the top-$t$ DP moves), and
the banded QUBO of width $k-1$ as a bridge to quantum neighborhood
search~\cite{windowed_own}.}

% ---------------------------------------------------------------------
\begin{thebibliography}{99}\small
\bibitem{ahuja2002} R.K.~Ahuja, \"O.~Ergun, J.B.~Orlin, A.P.~Punnen.
A survey of very large-scale neighborhood search techniques.
\emph{Discrete Appl.\ Math.} 123:75--102, 2002.
\bibitem{balzereit2024} \todo{verify authors/venue} Balzereit et al.
Parameterized local search for scheduling. arXiv:2409.00771, 2024.
\bibitem{bozejko2007} W.~Bo\.zejko, M.~Wodecki. On the theoretical
properties of swap multimoves. \emph{Oper.\ Res.\ Lett.} 35(2):227--231, 2007.
\bibitem{bozejko2026sr} W.~Bo\.zejko, M.~Uchro\'nski, M.~Wodecki. The new
exponential Fibonacci neighborhood as QUBO generated on a D-Wave quantum
computing environment for the permutational scheduling problem.
\emph{Sci.\ Rep.}, 2026. doi:10.1038/s41598-026-66284-9.
\bibitem{congram2002} R.K.~Congram, C.N.~Potts, S.L.~van de Velde. An
iterated dynasearch algorithm for the single-machine total weighted
tardiness scheduling problem. \emph{INFORMS J.\ Comput.} 14(1):52--67, 2002.
\bibitem{ergun2006} \todo{verify} \"O.~Ergun, J.B.~Orlin. Fast
neighborhood search for the single machine total weighted tardiness
problem. \emph{Oper.\ Res.\ Lett.} 34(1):41--45, 2006.
\bibitem{fibonacci_own} \todo{own work, third person: Fibonacci neighborhood, ICAISC 2026 / BPASTS}
\bibitem{garey1976} M.R.~Garey, D.S.~Johnson, R.~Sethi. The complexity of
flowshop and jobshop scheduling. \emph{Math.\ Oper.\ Res.} 1(2):117--129, 1976.
\bibitem{grabowski2004} J.~Grabowski, M.~Wodecki. A very fast tabu search
algorithm for the permutation flow shop problem with makespan criterion.
\emph{Comput.\ Oper.\ Res.} 31(11):1891--1909, 2004.
\bibitem{grosso2004} \todo{verify} A.~Grosso, F.~Della Croce, R.~Tadei. An
enhanced dynasearch neighborhood for the single-machine total weighted
tardiness scheduling problem. \emph{Oper.\ Res.\ Lett.} 32(1):68--72, 2004.
\bibitem{motzkin_own} \todo{own work, third person: Motzkin neighborhood, SOCO 2026, CCIS 3046, pp.~144--155}
\bibitem{nawaz1983} M.~Nawaz, E.E.~Enscore, I.~Ham. A heuristic algorithm
for the $m$-machine, $n$-job flow-shop sequencing problem. \emph{Omega}
11(1):91--95, 1983.
\bibitem{nowicki1996} E.~Nowicki, C.~Smutnicki. A fast tabu search
algorithm for the permutation flow-shop problem. \emph{Eur.\ J.\ Oper.\ Res.}
91(1):160--175, 1996.
\bibitem{ruiz2007} R.~Ruiz, T.~St\"utzle. A simple and effective iterated
greedy algorithm for the permutation flowshop scheduling problem.
\emph{Eur.\ J.\ Oper.\ Res.} 177(3):2033--2049, 2007.
\bibitem{smutnicki1998} C.~Smutnicki. \emph{Algorytmy szeregowania}.
Akademicka Oficyna Wydawnicza EXIT, Warszawa, 1998. \todo{check year/edition}
\bibitem{taillard1993} E.~Taillard. Benchmarks for basic scheduling
problems. \emph{Eur.\ J.\ Oper.\ Res.} 64(2):278--285, 1993.
\bibitem{windowed_own} \todo{own work, third person: windowed QUBO, under review}
\end{thebibliography}
\end{document}
